\section{Безопасность и проверяемость}
\label{sec:security}

При оффлоаде prefill поставщик присылает клиенту память модели, и через неё можно управлять генерацией. Ганеш и соавторы подменяли сегмент KV-кэша моделей Qwen3 кэшем на другую тему и меняли ход генерации,
не трогая видимый промпт \cite{ganesh2025historyswap}. Похожие атаки описаны для переиспользуемых кэшей рассуждений
\cite{zhu2026r2ecap} и для кэшей, которые переносят между позициями \cite{zhang2026hijackkv}. Если память модели
отравлена, модель сама решает сделать нужный атакующему вызов инструмента. То, что вызов исполняет клиент, от этого
не защищает, так же как не защищает и в облаке.

\textbf{Что меняет гибрид.} В облачном инференсе клиент проверяет сервер только с его помощью: проверяемый инференс
требует, чтобы поставщик сам публиковал хэши своих вычислений \cite{ong2025toploc}, а повторить вычисление клиенту
нечем. В гибридной схеме у клиента та же модель, что у сервера, и он может проверить присланную дельту,
посчитав её сам. Проверка устроена так.
\begin{itemize}
\item Журнал фиксирует историю. Дельта несёт идентификаторы токенов и хэш префикса; если они расходятся с журналом
клиента, клиент её отклоняет, поэтому поставщик не может переписать историю на уровне токенов.
\item Клиент выборочно пересчитывает дельты. Перед вставкой дельты он сохраняет своё рекуррентное состояние, 157~МБ,
а в простое пересчитывает случайно выбранную дельту от своего кэша и сравнивает результат с тем, что прислал сервер.
Начинать стоит с рекуррентного состояния: это одна сводка всей истории для большинства слоёв, она не привязана к
позициям и заменяется целиком при каждом оффлоаде, поэтому для атакующего это самый удобный рычаг.
\item Проверка может идти по частям. Сервер вместе с дельтой присылает хэши промежуточных состояний каждого слоя на
границах кусков по несколько десятков токенов. Клиент выбирает случайные пары из слоя и куска, запрашивает их
входы, пересчитывает только этот слой на этом куске и сверяет с хэшем. Чтобы подделать итоговое состояние, сервер
должен подделать хотя бы один переход в этой цепочке, и случайные проверки находят его с вероятностью, которая
растёт с их числом. Близкая схема с хэшами промежуточных активаций описана как проверяемый инференс
\cite{ong2025toploc}.
\end{itemize}
Полный пересчёт стоит столько же, сколько локальное чтение этой порции: тысяча токенов на глубине 170 тысяч занимает
на M3 Max по кривой Q4\_K\_M на Metal больше полуминуты. Если проверять 5\% дельт, клиент повторяет около 5\% работы, которую отдал серверу, и делает
это, когда простаивает. Проверка по частям дешевле полного пересчёта; во сколько, зависит от размера кусков и числа проверок, и этот протокол
мы предлагаем, но не измеряли. Выборочный пересчёт находит подмену после ответа и не отменяет вызов инструмента,
который модель уже сделала; покрытие атак такой проверкой мы не оценивали.

\textbf{Точная сверка и сверка с допуском.} Если сервер и клиент считают одни и те же веса одной и той же арифметикой
и с одной разбивкой на пакеты, при детерминированных ядрах вычислений их числа совпадают бит в бит; на
CPU-прототипе так и вышло (приложение~\ref{app:prototype}). Тогда любое отличие означает подмену, и сверка точная.
Если сервер считает на CUDA, а клиент на Metal, или кванты у сторон разные, числа честного сервера расходятся с пересчётом клиента; насколько, для пар CUDA и Metal и для разнородных квантов
нужно измерить. Тогда клиент задаёт порог, например наибольшее допустимое расхождение
распределений следующего токена, и считает всё, что ниже порога, честной погрешностью. Это сверка с допуском. Она
ловит грубую подмену, а опубликованные атаки грубые: они заменяют целые сегменты кэша, и расхождение получается на
порядки больше погрешности. Подмену, которую атакующий подогнал под порог, сверка с допуском не отличит от
погрешности. Отсюда цена разнородных квантов из раздела~\ref{sec:architecture}: они делают генерацию клиента быстрее,
а проверку поставщика слабее. Лоскутный режим оставляет поставщику меньшую долю кэша внимания, чем ресинк, но рекуррентное
состояние в обоих режимах серверное. Поэтому с недоверенным поставщиком клиент работает в лоскутном режиме и в
первую очередь выборочно пересчитывает рекуррентное состояние. Выборочный пересчёт обнаруживает
подмену. Если клиент сверяет память после того, как агент на ней сгенерировал ответ и выполнил действие, подмена
обнаружится уже после действия. Проверка нужной части памяти до генерации снижает риск, что агент будет работать на
подменённой памяти. Гарантий проверка не даёт: выборочный пересчёт может пропустить подмену вне проверенных участков, а
сверка с допуском пропускает подмену в пределах допуска. Насколько проверка снижает риск, зависит от покрытия проверки и
модели угроз. Платит клиент за проверку частью выигрыша во времени.

\textbf{Уровни поставщиков.} Из этого вырастают два уровня услуги. Проверяемый prefill работает на тех же весах и
детерминированной арифметике, что и клиент, и позволяет точную сверку; он медленнее и дороже. Оптимизированный
prefill использует самые быстрые движки и форматы сервера, например NVIDIA с 4-битными весами, при клиенте на MLX, и
стоит дешевле ценой проверяемости. Пользователь с чувствительными задачами выберет первый; так может возникнуть рынок проверяемого prefill, которого в
облачном инференсе нет.

\textbf{Чего схема не решает.} Поставщик prefill видит все токены сессии, то есть код и данные пользователя, и
конфиденциальность схема не обеспечивает: её дают только доверие к поставщику, конфиденциальные вычисления или
выбор нечувствительных задач. Полная модель угроз против отравления памяти модели и её проверки выходит за рамки
этой работы. Мы показываем, что гибридная схема даёт клиенту механизмы проверки, которых у облачного инференса нет,
и где эти механизмы перестают работать.

\textbf{Вывод.} Гибрид даёт клиенту способ обнаружить подмену присланной памяти модели, которого нет в облаке, но
конфиденциальности не даёт. Работа с поставщиком, которому клиент не доверяет, требует как минимум тех же весов,
детерминированных вычислений и выборочного пересчёта, в первую очередь рекуррентного состояния; достаточно ли
этого, покажет модель угроз.
